The model is solved under both tax regimes and calibrated to a joint-taxation baseline.
The counterfactual switches the regime to individual taxation on a revenue-neutral basis.
I compare the two steady states along three dimensions: (i) the gender gap in hours worked at each age, (ii) the lifetime human capital trajectories of both spouses, and (iii) household welfare measured by the consumption-equivalent variation.

The central prediction is that individual taxation raises the secondary earner's labor supply directly through the lower effective marginal rate, and indirectly through the human capital channel: more hours today raise future wages, which further increase labor supply in subsequent periods.
The dynamic amplification separates this framework from a standard static tax analysis and is the key contribution of the proposed model.

As an additional exercise, the transition path from joint to individual taxation can be simulated to quantify adjustment costs and the speed at which the gender gap closes following a reform.
